\chapter{Planificación del primer incremento}
\section{Referencia histórica y propuesta actual}
La referencia histórica proporcionada por el usuario sitúa HU-01 a HU-04 en Sprint 1. No se han consultado las estimaciones originales de Jira y no se modifican. HU-16 se identifica como dependencia técnica y funcional del acceso protegido, sin afirmar que estuviera incluida históricamente.

La planificación propuesta se conserva en \path{docs/SPRINT_1.md}. La duración es dos semanas, pero no se inventan fechas de eventos, horas disponibles o velocidad del equipo. Antes de comprometer todo el alcance se deben contrastar disponibilidad, dificultad y criterios con los cuatro integrantes.

\section{Objetivo verificable}
Una cuenta administradora de demostración puede iniciar sesión, registrar una cuenta con un rol permitido, crear una capacitación, editarla y volver a consultar los datos persistidos. Las solicitudes sin acceso, los datos inválidos y los conflictos reciben errores comprensibles; el navegador no simula un guardado local.

\begin{table}[htbp]
\centering
\caption{Paquetes de trabajo del primer incremento}
\begin{tabular}{P{2cm}P{7cm}P{3cm}}
\toprule Referencia & Resultado esperado & Verificación \\ \midrule
HU-16 & Login, sesión vigente y logout & API y revisión de navegador \\
HU-01 & Alta, consulta y edición de usuarios & Validez, duplicado y acceso \\
HU-02 & Roles, asignaciones y permisos & Privilegios y administrador protegido \\
HU-03 & Crear capacitación persistente & Caso válido y límites \\
HU-04 & Editar sin violar integridad & Consulta posterior y estados \\
Técnica & Migraciones, semilla, documentación & Arranque reproducible \\
\bottomrule
\end{tabular}
\end{table}

\section{Entrega mínima y capacidad}
Si no existe capacidad para terminar todo en el Sprint, la entrega mínima debe contener login, un usuario autorizado y creación/consulta persistente de capacitación, con aceptación y pruebas. Edición completa de permisos u otras tareas deben registrarse pendientes, sin declarar terminado un incremento que no cumple su objetivo.

La documentación puede elaborarse en paralelo con límites de archivos claros, pero una revisión de código por otro integrante sigue siendo necesaria. Las tareas compartidas deben integrarse de forma consciente; las contribuciones de IA no representan automáticamente trabajo de cada estudiante.

\section{Evolución propuesta}
V2 revisará el diseño real y ampliará instructores y participantes cuando sus dependencias estén listas. V3 podrá incorporar sesiones, asistencia, historial y evaluaciones. V4 priorizará integración, seguridad, experiencia y funcionalidades restantes acordadas. Ninguna de estas metas se etiqueta como versión ya construida o probada.
